\documentclass[aspectratio=169,10pt]{beamer}
\usetheme{metropolis}
\metroset{numbering=fraction,progressbar=frametitle,block=fill}
\setbeamertemplate{navigation symbols}{}
\usepackage{amsmath,amssymb,mathtools,booktabs,array,tabularx,ragged2e}
\usepackage{xeCJK}
\setCJKmainfont[BoldFont=Hiragino Sans GB W6]{Hiragino Sans GB W3}
\setCJKsansfont[BoldFont=Hiragino Sans GB W6]{Hiragino Sans GB W3}
\setmainfont{Helvetica Neue}
\setsansfont{Helvetica Neue}
\usepackage{newunicodechar}
\newfontfamily\symfont{STIX Two Math}
\newunicodechar{→}{{\symfont →}}\newunicodechar{⇒}{{\symfont ⇒}}\newunicodechar{⇔}{{\symfont ⇔}}\newunicodechar{↔}{{\symfont ↔}}\newunicodechar{≡}{{\symfont ≡}}\newunicodechar{≤}{{\symfont ≤}}\newunicodechar{≥}{{\symfont ≥}}\newunicodechar{≠}{{\symfont ≠}}\newunicodechar{∈}{{\symfont ∈}}\newunicodechar{∞}{{\symfont ∞}}\newunicodechar{×}{{\symfont ×}}\newunicodechar{·}{{\symfont ·}}\newunicodechar{①}{{\symfont ①}}\newunicodechar{②}{{\symfont ②}}\newunicodechar{③}{{\symfont ③}}\newunicodechar{④}{{\symfont ④}}
\definecolor{positive}{HTML}{0173B2}
\definecolor{negative}{HTML}{DE8F05}
\definecolor{emphasis}{HTML}{029E73}
\definecolor{neutral}{gray}{0.55}
\newcommand{\pos}[1]{\textcolor{positive}{#1}}
\newcommand{\con}[1]{\textcolor{negative}{#1}}
\newcommand{\HL}[1]{\textcolor{emphasis}{#1}}
\setbeamerfont{frametitle}{size=\large}
\setbeamerfont{title}{size=\LARGE}
\setlength{\leftmargini}{1.2em}
\newcolumntype{Y}{>{\RaggedRight\arraybackslash}X}
\title{GFlowNet × 最优传输}
\subtitle{从流守恒到 Kantorovich 计划：18 篇论文、两条趋势线、一个可开工的课题}
\author{awesome\_Gflow\_OT ???}
\institute{github.com/asimfish/awesome\_Gflow\_OT}
\date{2026-09}
\begin{document}
\maketitle
\begin{frame}{这份汇报只想说清五件事}\framesubtitle{一页结论}
\begin{itemize}\small
\item \textbf{GFlowNet 的内部流从来不唯一}：奖励只钉住终止边，非终止边上的 $P_B$ 自由（T02 Prop. 18）；有环时再加一个环空间（T19 Prop. 5）
\item \textbf{最小总流是一个合法的选择原则}，因为 $\sum_s F(s)=Z\,\mathbb E[n_\tau]$（T36）：总流最小 = 期望轨迹最短
\item \textbf{固定源分布后，最小总流 = 图最短路代价下的 Kantorovich OT}（O08 Thm. 3.2），对偶势 + 互补松弛给出证书（Thm. 3.3）
\item \textbf{等价是经典的（Beckmann / 最小费用流），接口是新的}：输出不是耦合矩阵，而是能在 $20!$ 规模隐式图上执行的局部策略
\item \textbf{该做什么}：\HL{balance 残差 → OT 误差界 + 对偶势证书}。撞车风险最低，OT 侧配件刚好齐了；摊销与熵正则两条路已经拥挤
\end{itemize}
\end{frame}

\section{问题}

\begin{frame}{奖励只钉住边界}\framesubtitle{GFlowNet 的两条约束}
\begin{equation*}\sum_{u}F(u\to s)=\sum_{v}F(s\to v),\qquad F(x\to s_f)=R(x)\quad\Longrightarrow\quad P_T(x)=\frac{R(x)}{Z}\end{equation*}
\begin{itemize}\small
\item 两条约束都是线性的（T02 Prop. 19 / Eq. (22)）：可行集是一个多面体
\item 奖励只出现在终止边上，它固定的是多面体的\textbf{边界}
\item T02 Prop. 18 第 3 条：Markovian flow 由终止流与非终止边上的 $P_B$ 唯一确定——**选一个 $P_B$ 就是选一个内部流**
\item 原文对「选哪个」只有一句（§2.6）：可以偏好更短的路径
\end{itemize}
\end{frame}

\begin{frame}{自由度更大，也更危险}\framesubtitle{有环之后}
\begin{columns}[T]
\begin{column}{0.48\textwidth}\textbf{\pos{T19 · AAAI 2024}}\vspace{2pt}
{\footnotesize\begin{itemize}
\item R-流集合 = 一个无环流 + 环空间 $H^1_+(G)$（Prop. 5）
\item 比值型损失（FM/DB/TB）把质量无限堆进环：flow explosion（Thm. 3）
\item 差值型损失稳定（Thm. 4）
\item 总流与期望长度只有不等式（Thm. 2）
\end{itemize}}\end{column}
\hfill
\begin{column}{0.48\textwidth}\textbf{\pos{T36 · ICML 2025}}\vspace{2pt}
{\footnotesize\begin{itemize}
\item 以 $P_B$ 为第一性对象，流 = 期望访问次数 × 终止流（Prop. 3.7 一一对应）
\item 不等式收紧为等式 $\sum_s F(s)=Z\,\mathbb E[n_\tau]$
\item 写出「最小总流」约束优化 Eq. (11)；近似解法 DB + 状态流正则
\item 最小值落在多面体边界，超出自身正性假设，只能 $\lambda$ 正则逼近
\end{itemize}}\end{column}
\end{columns}
\vfill\hrule\vspace{2pt}{\footnotesize\textcolor{neutral}{易误读：T19「极限流无环」≠「解唯一」。DAG 上环空间为零时 R-流照样不唯一；真正选出一个点的是 T36 的线性目标。}}
\end{frame}

\begin{frame}{每一环缺什么、下一环补什么}\framesubtitle{主线逻辑链}
\begin{columns}[T]
\begin{column}{0.182\textwidth}\begin{block}{T02}\footnotesize 流守恒 ⇒ $P_T\propto R$；$P_B$ 自由\\[3pt]\textcolor{negative}{缺：选哪个内部流}\end{block}\end{column}
\begin{column}{0.182\textwidth}\begin{block}{T19}\footnotesize 有环流理论；环空间；flow explosion\\[3pt]\textcolor{negative}{缺：总流 vs 长度只有 ≤}\end{block}\end{column}
\begin{column}{0.182\textwidth}\begin{block}{T36}\footnotesize $\sum F(s)=Z\,\mathbb E[n_\tau]$；最小流优化\\[3pt]\textcolor{negative}{缺：最小流是什么}\end{block}\end{column}
\begin{column}{0.182\textwidth}\begin{block}{O07}\footnotesize 期望长度最小 ⇔ 只走最短路（Thm. 3.4）\\[3pt]\textcolor{negative}{缺：单源、质量怎么分}\end{block}\end{column}
\begin{column}{0.182\textwidth}\begin{block}{O08}\footnotesize 固定源分布 ⇒ Kantorovich OT（Thm. 3.2）\\[3pt]\textcolor{negative}{缺：训练与 LP 之间无界}\end{block}\end{column}
\end{columns}
\vfill\hrule\vspace{2pt}{\footnotesize\textcolor{neutral}{转折点是 T36 的等式：一个流空间上的线性泛函（总流）和一个行为学量（期望步数）被钉在一起，此后 O07/O08 是线性规划课的两道习题。}}
\end{frame}

\section{核心定理}

\begin{frame}{最小总流 = Kantorovich 最优传输}\framesubtitle{O08 Theorem 3.2}
\begin{equation*}\min_{\mathcal F\ge0}\ \sum_{e\in E^\circ}\mathcal F(e)\ \ \text{s.t.}\ \operatorname{div}\mathcal F=L-R\quad=\quad\min_{\Pi\in\Gamma(L,R)}\ \sum_{u,x}d_G(u,x)\,\Pi(u,x)\end{equation*}
\begin{itemize}\small
\item 四步证明：双线性 DB 约束 → 消策略成 LP（App. A.1）→ 约简成 divergence 约束（Eq. (10)→(11)）→ 两侧夹逼
\item $\le$：取最优耦合 $\Pi^\star$，沿每对 $(u,x)$ 的一条最短路铺流，可行且代价相等
\item $\ge$：取最优流 $\mathcal F^\star$，由它诱导策略与轨迹分布，端点分布 $\Pi^\star_{u,x}=\sum_{\tau:u\rightsquigarrow x}\mathbb P^\star(\tau)$ 是可行耦合
\item Thm. 3.3（对偶）：$\pi^\star_x=d(x)$，互补松弛 $\mathcal F^\star(s\to s')(\pi^\star_{s'}-1-\pi^\star_s)=0$——最优流只走最短路子图
\end{itemize}
\end{frame}

\begin{frame}{学到的策略同时满足边缘与 OT 最优代价}\framesubtitle{O08 Table 1 · Hypergrid（数字照抄原文）}
{\footnotesize\setlength{\tabcolsep}{3pt}\begin{tabular}{@{}>{\RaggedRight\arraybackslash}p{0.147\textwidth}>{\RaggedRight\arraybackslash}p{0.086\textwidth}>{\RaggedRight\arraybackslash}p{0.184\textwidth}>{\RaggedRight\arraybackslash}p{0.184\textwidth}>{\RaggedRight\arraybackslash}p{0.184\textwidth}>{\RaggedRight\arraybackslash}p{0.184\textwidth}@{}}\toprule
\textbf{源分布 L} & \textbf{H} & \textbf{TV（学到）} & \textbf{TV*（完美采样器）} & \textbf{E|τ|（学到）} & \textbf{OT*} \\\midrule
Ball & 10 & 0.024 & 0.024 & 3.990 & 3.997 \\\addlinespace[2pt]
Ball & 15 & 0.036 & 0.033 & 6.325 & 6.303 \\\addlinespace[2pt]
Ball & 20 & 0.037 & 0.040 & 8.326 & 8.325 \\\addlinespace[2pt]
Moon & 10 & 0.022 & 0.024 & 4.352 & 4.351 \\\addlinespace[2pt]
Moon & 15 & 0.023 & 0.033 & 6.907 & 6.868 \\\addlinespace[2pt]
Moon & 20 & 0.032 & 0.040 & 9.001 & 9.059 \\\bottomrule\end{tabular}}
\vfill\hrule\vspace{2pt}{\footnotesize\textcolor{neutral}{Figure 1：H=10 Moon 上精确 LP 解的边流与 Kantorovich 最优耦合的 transport cost 都是 4.351——同一个数，一边是边流、一边是耦合。排列 S\_20（Table 2）无 OT* 参考，只有边缘诊断。}}
\end{frame}

\begin{frame}{Theorem 3.2 在什么条件下成立}\framesubtitle{O08 的八条前提}
\begin{columns}[T]
\begin{column}{0.235\textwidth}\begin{block}{图结构}\footnotesize 有限有向图，\textbf{允许有环}（无环分层图上路径等长，最小流退化为常数）；$s_0$ 只出到 $U$，$s_f$ 只从 $\mathcal X$ 入；全对可达\end{block}\end{column}
\begin{column}{0.235\textwidth}\begin{block}{边缘与代价}\footnotesize 两个边缘都归一化：$\sum L=\sum R=1$，$Z$ 已知；代价是\textbf{单位边长}的最短路跳数，带权图不在覆盖范围\end{block}\end{column}
\begin{column}{0.235\textwidth}\begin{block}{吸收性}\footnotesize 反向链是吸收链 / $\mathbb E[n_\tau]<\infty$（T36 Lemma 3.4；O07 Assumption 3.1）；正文未重述但 Eq. (3) 良定义依赖它\end{block}\end{column}
\begin{column}{0.235\textwidth}\begin{block}{精确 vs 训练}\footnotesize 定理是关于 \textbf{LP 最优解} 的；神经训练用软罚 $\lambda$（Eq. (20)），训练完的模型是否 OT 最优靠数值接近程度支撑，不由定理保证\end{block}\end{column}
\end{columns}
\vfill\hrule\vspace{2pt}{\footnotesize\textcolor{neutral}{适用范围一句话：状态空间无法枚举、代价矩阵无法构造、但局部合法动作明确且允许回退的组合图上的均衡 OT。不适合显式小图（network simplex 更快更准且自带证书）。}}
\end{frame}

\begin{frame}{等价是经典的，接口是新的}\framesubtitle{新意到底在哪}
\begin{columns}[T]
\begin{column}{0.48\textwidth}\textbf{\pos{不是新的}}\vspace{2pt}
{\footnotesize\begin{itemize}
\item 图上最短路代价的 OT ≡ 最小费用流 ≡ 离散 Beckmann（Beckmann 1952；O02 Essid \& Solomon 2018；O01 Prop. 6.23 / 14.9）
\item 流分解定理、最优流无环、LP 顶点解稀疏（O01 Prop. 3.8）
\item 「对偶势 = 最短路距离」是最小费用流对偶的教科书结论
\end{itemize}}\end{column}
\hfill
\begin{column}{0.48\textwidth}\textbf{\pos{新的}}\vspace{2pt}
{\footnotesize\begin{itemize}
\item 解是只依赖当前状态的局部策略 $P_F(s'\mid s)$，不是 $|U|\times|\mathcal X|$ 矩阵
\item 可在无法枚举、无法构造代价矩阵的组合图上执行（$S_{20}$，$2.4\times10^{18}$ 状态）
\item GFlowNet 的现成训练基础设施（T36 代码、gfnx / torchgfn）直接复用
\item \con{代价}：为了成为 LP，放弃了「只需未归一化奖励」这一招牌能力
\end{itemize}}\end{column}
\end{columns}
\end{frame}

\section{OT 侧}

\begin{frame}{OT 概念 ↔ GFlowNet 概念}\framesubtitle{从 O01（Peyré 讲义）提炼}
{\scriptsize\setlength{\tabcolsep}{3pt}\begin{tabular}{@{}>{\RaggedRight\arraybackslash}p{0.286\textwidth}>{\RaggedRight\arraybackslash}p{0.184\textwidth}>{\RaggedRight\arraybackslash}p{0.212\textwidth}>{\RaggedRight\arraybackslash}p{0.288\textwidth}@{}}\toprule
\textbf{OT} & \textbf{GFlowNet} & \textbf{关系} & \textbf{出处} \\\midrule
边流 $m_e$，目标 $\sum_e\ell_e|m_e|$ & 边流 $F(s\to t)$，总流 & \pos{严格等价} & O01 Prop. 6.23；O02 Eq. (1)⇔(3) \\\addlinespace[2pt]
守恒律 $\operatorname{div}_G m=r$ & flow matching & \pos{严格等价} & O01 Prop. 6.23 \\\addlinespace[2pt]
代价 = 图最短路 & 轨迹长度 / 期望步数 & 条件等价（路径作用 = 边长和） & O01 Prop. 14.9 \\\addlinespace[2pt]
边缘约束 & 奖励匹配 + 第一步边流 = L & 条件等价（$Z=1$） & O08 Assumption 3.1 \\\addlinespace[2pt]
互补松弛 & 最优流只在紧边为正 & \pos{严格}（LP 形式） & O01 Prop. 5.3；O08 Thm. 3.3 \\\addlinespace[2pt]
对偶势 $f_i$ & 状态流 $\log F(s)$ & 类比（不等式 vs 等式） & O01 Def. 5.1 \\\addlinespace[2pt]
熵正则 $\varepsilon\mathrm{KL}$ & 路径熵 & 类比（端点熵 vs 路径熵） & O01 Def. 8.2 \\\addlinespace[2pt]
Schrödinger 桥参考律 & 固定的 $P_B$ & 条件等价 & O01 Prop. 14.11；O04 投影同构 \\\addlinespace[2pt]
Benamou–Brenier / JKO & 训练动力学 & \con{仅氛围类比} & O01 Thm. 14.5 \\\bottomrule\end{tabular}}
\vfill\hrule\vspace{2pt}{\footnotesize\textcolor{neutral}{易误读：单源 DAG 上「学到 OT 计划」是空话——$\mathcal U(\delta_{s_0},\beta)$ 是单点集（O01 Remark 3.2）。}}
\end{frame}

\begin{frame}{四篇神经 OT / SB：谁威胁、谁互补}\framesubtitle{相邻工作}
\begin{columns}[T]
\begin{column}{0.235\textwidth}\begin{block}{O05 UNOT · ICML 2025 · \con{威胁最大}}\footnotesize 跨数据集/跨分辨率摊销熵 OT 对偶势，非训练集相对误差 1.3–2.8\%（Table 4，Meta OT 11–31\%）。它摊销的势正是 O08 Thm. 3.3 的 $\pi_x$ 连续版；移植到图最短路代价是周末量级。GFN 剩下的生态位：换图即换代价、状态空间存不下代价矩阵\end{block}\end{column}
\begin{column}{0.235\textwidth}\begin{block}{O04 α-DSBM · NeurIPS 2024 Spotlight · \pos{互补最大}}\footnotesize $\mathrm{proj}_{\mathcal R}/\mathrm{proj}_{\mathcal M}$ 与「固定 $P_B$」/「由边流反解 $P_F$」逐条同构；Eq. (26) + Lemma D.2 是「图上 SB = 熵正则 min-flow GFN」的现成证明模板。易误读：$\alpha$ 是耦合空间插值系数，不是学习率\end{block}\end{column}
\begin{column}{0.235\textwidth}\begin{block}{O03 GeONet · UAI 2024}\footnotesize 神经算子学 Wasserstein 测地线，连续分布 + 动态 OT；与 GFN–OT 的互补点是连续 geodesic 而非离散构造路径\end{block}\end{column}
\begin{column}{0.235\textwidth}\begin{block}{O06 · AISTATS 2025}\footnotesize 用连续正规化流的速度场解高维 OT。名字里的 flow neural network \textbf{不是 GFlowNet}\end{block}\end{column}
\end{columns}
\end{frame}

\section{竞争格局}

\begin{frame}{五篇的身份牌与分界线}\framesubtitle{COMPETITOR\_MATRIX（发表状态四路核实）}
{\scriptsize\setlength{\tabcolsep}{3pt}\begin{tabular}{@{}>{\RaggedRight\arraybackslash}p{0.086\textwidth}>{\RaggedRight\arraybackslash}p{0.237\textwidth}>{\RaggedRight\arraybackslash}p{0.161\textwidth}>{\RaggedRight\arraybackslash}p{0.139\textwidth}>{\RaggedRight\arraybackslash}p{0.205\textwidth}>{\RaggedRight\arraybackslash}p{0.142\textwidth}@{}}\toprule
\textbf{} & \textbf{O08 GFN-OT} & \textbf{O07 GFN-SP} & \textbf{C01 ULOT} & \textbf{C02 GSBoG} & \textbf{C03 DDSBM} \\\midrule
发表 & \con{ICML 2026 Workshop} & \con{ICML 2026 Workshop} & NeurIPS 2025 主会 & \textbf{ICML 2026 主会} & ICLR 2025 主会 \\\addlinespace[2pt]
图是什么 & 场所（隐式） & 场所（隐式） & 对象（显式） & 场所（显式） & 对象 \\\addlinespace[2pt]
cost / 正则 & 线性最短路，$\varepsilon=0$ & 最短路 & 二次 FUGW & KL 到参考 CTMC，$\varepsilon>0$ & KL，$\varepsilon>0$ \\\addlinespace[2pt]
输出 & \textbf{policy} & policy & plan 矩阵 & policy & 生成模型 \\\addlinespace[2pt]
可摊销 & 否 & 否 & \textbf{是} & 否 & 否 \\\addlinespace[2pt]
误差保证 & \textbf{LP 对偶 + 互补松弛} & Thm. 3.4 & 零定理 & 无误差界 & IMF 收敛无速率 \\\addlinespace[2pt]
最大规模 & $S_{20}$ & 3×3×3 Rubik & $n\le10^4$ & **$10^6$ 节点** & ZINC250K \\\addlinespace[2pt]
代码 & 未见 & 开源 & 开源 & \textbf{未公开} & 开源 \\\bottomrule\end{tabular}}
\vfill\hrule\vspace{2pt}{\footnotesize\textcolor{neutral}{GSBoG 与 O08 同象限、卖点句几乎逐字相同（可执行策略而非静态耦合），却在主会且规模高四个数量级。反制点只有三个：隐式图、不定时域、误差证书。}}
\end{frame}

\section{2025–2026 趋势}

\begin{frame}{两侧同时成熟，却没有人接线}\framesubtitle{趋势判断（每条 ≥2 篇支撑）}
\begin{columns}[T]
\begin{column}{0.48\textwidth}\textbf{\pos{GFlowNet 侧}}\vspace{2pt}
{\footnotesize\begin{itemize}
\item 非无环理论进主会（T36、Ergodic GF），三个应用出口（最短路、OT、MCMC 终止）全停 Workshop
\item \textbf{误差证书成为竞争点}：Stable GFlowNets 的 loss→TV 界、Evaluation Balance 把残差当评估器
\item 训练目标进入「族」时代：$f$-TB（ICML 2026）、$\alpha$-GFN、Hybrid-Balance（NeurIPS 2025）
\item 策略梯度回流：PPO for GFN、自然梯度、RTB ≡ Trust-PCL
\item $Z$ 的两种命运：LLM 场景里被消灭或被重用；O08 里根本不出现
\end{itemize}}\end{column}
\hfill
\begin{column}{0.48\textwidth}\textbf{\pos{最优传输侧}}\vspace{2pt}
{\footnotesize\begin{itemize}
\item \textbf{图上 OT 的理论重心在「势」}：对偶预测加速最小费用流（AAAI 2026）、Sinkhorn 势统计率（Belomestny）、QOT 的 PL 不等式
\item 离散 SB 完成闭环：有解析解基准 + 收敛率，新进入者无定义问题的红利
\item 摊销 OT 拥挤度继续上升（切片势摊销、min-sliced 计划）
\item 二次正则是唯一未被占满的方向；「非唯一时选哪个」有独立文献（熵选择原理）
\end{itemize}}\end{column}
\end{columns}
\end{frame}

\begin{frame}{窗口状态与最终评级}\framesubtitle{四个候选课题}
{\footnotesize\setlength{\tabcolsep}{3pt}\begin{tabular}{@{}>{\RaggedRight\arraybackslash}p{0.238\textwidth}>{\RaggedRight\arraybackslash}p{0.141\textwidth}>{\RaggedRight\arraybackslash}p{0.591\textwidth}@{}}\toprule
\textbf{课题} & \textbf{评级} & \textbf{依据} \\\midrule
\textbf{① Balance 残差 → OT 误差界 + 对偶势证书} & \pos{做} & GFN 侧有 loss→TV 界；OT 侧误差理论全在势上；O08 Thm. 3.3 的互补松弛给逐边证书；\textbf{没有一篇把近似可行流的残差映射到 cost gap} \\\addlinespace[2pt]
② 条件 GFN 摊销一族图上 OT & \con{不做}（降为对照） & UNOT、ULOT、切片势摊销、min-sliced 四篇以上；唯一护城河是隐式图，但「一族隐式图」怎么定义没人说清 \\\addlinespace[2pt]
③ 熵正则 GFN–OT / 图上 SB & \con{不做}（作①的 ε>0 推广） & GSBoG 占同象限主会位；Sampling Decisions 已把 GFN 流函数写成单侧 Schrödinger 传输；离散 SB 有基准与收敛率 \\\addlinespace[2pt]
④ GFN proposal + 经典 OT 修正 & 改形态 & 已被「对偶预测 → ε-relaxation」具体化（AAAI 2026）；可作①的应用：把学到的状态流当对偶预测喂经典求解器 \\\bottomrule\end{tabular}}
\vfill\hrule\vspace{2pt}{\footnotesize\textcolor{neutral}{如果只做一件事，做①；②③④作为①的基线或推广出现，而不是独立论文。}}
\end{frame}

\begin{frame}{决定性实验}\framesubtitle{可立即开工}
\begin{columns}[T]
\begin{column}{0.48\textwidth}\textbf{\pos{环境 · 方法 · 基线}}\vspace{2pt}
{\footnotesize\begin{itemize}
\item Hypergrid $H\in\{10,15,20\}$（精确 LP / network simplex 可得）→ $S_n$ Cayley 图 $n\in\{4,8\}$（精确）→ $n=20$（只报证书）→ 带权变体
\item 方法：min-flow GFN（T36 代码）× \{DB 正则, PPO 训练器, $f$-TB 损失\}
\item 基线：network simplex、Sinkhorn（含解析解基准校验）、UNOT 移植版、GSBoG 复现/α-DSBM 离散化替身
\end{itemize}}\end{column}
\hfill
\begin{column}{0.48\textwidth}\textbf{\pos{指标 · 判据}}\vspace{2pt}
{\footnotesize\begin{itemize}
\item 边缘：终止分布 TV、源边缘违反量；代价：$\mathbb E|\tau|$ vs OT*、$\|\Pi_\theta-\Pi^\star\|_1$
\item \textbf{证书}：由状态流构造对偶可行 $\pi$，报 primal–dual gap；核心图：x = balance 残差，y = cost gap
\item 预期：显式小图精确求解器全面占优；隐式图上 GFN 是唯一能给证书的方法
\item \con{失败判据}：$n=20$ 上 gap 与残差 Spearman < 0.5 → 只认证边缘误差；带权变体证书失效 → 单位边长假设是本质的
\end{itemize}}\end{column}
\end{columns}
\end{frame}

\section{仓库与方法论}

\begin{frame}{仓库里有什么}\framesubtitle{交付物}
\begin{columns}[T]
\begin{column}{0.235\textwidth}\begin{block}{18 篇深度解读}\footnotesize 严格 8 节模板；每个定理、数字带原文出处；末尾「编者注」记录自主决定的歧义。竞品矩阵、两份趋势扫描、INSIGHTS 综合\end{block}\end{column}
\begin{column}{0.235\textwidth}\begin{block}{17 篇保版式中译 PDF}\footnotesize SuperTranslate：公式/图表/引用冻结，正文按原坐标回填；每篇附对象级 QA（inspect.json）。O01 讲义因 API 余额未译\end{block}\end{column}
\begin{column}{0.235\textwidth}\begin{block}{汇总报告 中/英}\footnotesize 由 \texttt{scripts/build\_report.py} 从解读报告拼装，pandoc + xelatex 编译：中文 91 页、英文 120 页（DeepSeek/Gemini 分块翻译，LaTeX 与引用逐字保留）\end{block}\end{column}
\begin{column}{0.235\textwidth}\begin{block}{数据驱动 README}\footnotesize \texttt{data/meta/*.json} 一篇一卡 → \texttt{src/generator.py} 生成中英 README（awesome-ml4co 模式）；2026 增补自 \texttt{candidates\_*.csv}\end{block}\end{column}
\end{columns}
\vfill\hrule\vspace{2pt}{\footnotesize\textcolor{neutral}{发表状态纪律：主会 / 期刊 / Workshop / 预印本分开标，Workshop 绝不写成主会；arXiv comment 滞后时以 OpenReview / dblp / proceedings 为准。}}
\end{frame}

\begin{frame}{一句话带走}\framesubtitle{结语}
\begin{quote}\large GFlowNet 唯一不可替代的产出是\textbf{校准的多样解分布}；在 OT 这条线上，它唯一不可替代的位置是\textbf{无法枚举、无法构造代价矩阵的组合图}。在这个位置上给出\textbf{带证书}的传输策略，是目前没人占的那个空位。\end{quote}\vspace{6pt}{\small\textcolor{neutral}{awesome\_Gflow\_OT · reports/INSIGHTS.md · 2026-09}}
\end{frame}
\begin{frame}[shrink=15]{参考文献（核心链）}
\tiny\setbeamertemplate{bibliography item}{}
\begin{thebibliography}{9}
\bibitem{O08} Maksimov, Morozov, Belomestny, Samsonov. \emph{Your GFlowNet Secretly Learns an Optimal Transport Plan}. ICML 2026 SPIGM Workshop. arXiv:2606.06272.
\bibitem{O07} Morozov, Maksimov, Tiapkin, Samsonov. \emph{Learning Shortest Paths with Generative Flow Networks}. ICML 2026 SPIGM Workshop. arXiv:2603.01786.
\bibitem{T36} Morozov et al. \emph{Revisiting Non-Acyclic GFlowNets in Discrete Environments}. ICML 2025. arXiv:2502.07735.
\bibitem{T19} Brunswic et al. \emph{A Theory of Non-Acyclic Generative Flow Networks}. AAAI 2024. arXiv:2312.15246.
\bibitem{T02} Bengio et al. \emph{GFlowNet Foundations}. JMLR 24(210), 2023. arXiv:2111.09266.
\bibitem{O02} Essid, Solomon. \emph{Quadratically Regularized Optimal Transport on Graphs}. SIAM J. Sci. Comput. 2018. arXiv:1704.08200.
\bibitem{O01} Peyr\'e. \emph{Optimal Transport for Machine Learners}. Lecture notes 2025. arXiv:2505.06589.
\bibitem{C02} \emph{Generalized Schr\"odinger Bridge on Graphs}. ICML 2026. arXiv:2602.04675. \quad \textbf{[C01]} ULOT, NeurIPS 2025, arXiv:2506.12025. \quad \textbf{[C03]} DDSBM, ICLR 2025, arXiv:2410.01500.
\end{thebibliography}
\end{frame}
\begin{frame}[standout]?? · ?????\end{frame}
\appendix
\begin{frame}{Backup Slides}\centering\Large 备份页\end{frame}
\begin{frame}{Backup · O08 Theorem 3.2 证明骨架}
\small
\begin{enumerate}
\item \textbf{消策略}：DB 约束 $F(s)P_F(s'\mid s)=F(s')P_B(s\mid s')$ 在决策变量里是双线性的；以边流 $\mathcal F(s\to s')$ 为变量后约束成为线性（Appendix A.1）。
\item \textbf{约简}：把状态流 $\mathcal F(s)=\sum_{s'}\mathcal F(s\to s')$ 代入，流守恒写成 $\operatorname{div}\mathcal F(s)=L(s)-R(s)$（Eq. (10)$\to$(11)）；总流目标去掉常数 $\sum_u L(u)=1$。
\item \textbf{GFlow$^\star\le$OT$^\star$}：取最优耦合 $\Pi^\star$，对每对 $(u,x)$ 任选一条最短路 $\tau_{u,x}$，令 $\mathcal F=\sum_{u,x}\Pi^\star_{u,x}\,\mathbb I[\cdot\in\tau_{u,x}]$；路径的 divergence 是 $\mathbb I[s=u]-\mathbb I[s=x]$，故可行，代价 $=\sum\Pi^\star d$。
\item \textbf{GFlow$^\star\ge$OT$^\star$}：取最优流 $\mathcal F^\star$，还原 $P_F^\star$ 与轨迹分布 $\mathbb P^\star$；端点分布 $\Pi^\star_{u,x}=\sum_{\tau:u\rightsquigarrow x}\mathbb P^\star(\tau)$ 是可行耦合，且每条轨迹长度 $\ge d(u,x)$，故总流 $\ge$ OT 代价。
\item \textbf{对偶}（Thm. 3.3）：$\max_\pi\sum_x R(x)\pi_x$ s.t. $\pi_{s_0}=0,\ \pi_{s'}-\pi_s\le1$；$\pi_s\le d(s)$，$d(\cdot)$ 对偶最优；互补松弛 $\Rightarrow$ 最优流支撑在最短路子图。
\end{enumerate}
\end{frame}
\begin{frame}{Backup · 中译 PDF 与 QA 口径}
\small
\begin{itemize}
\item 引擎：SuperTranslate \texttt{pdf\_zh\_translator}，参数 \texttt{--preserve-graphics-text --skip-overflow}，DeepSeek 后端。
\item 不重排页面：公式、图表、引用先冻结，正文按原坐标回填；每页 \texttt{inspect} 对象级比对。
\item 17/18 完成；O01（Peyr\'e 讲义，百余页）因翻译 API 余额耗尽未译。
\item 已知 issue 类型：\texttt{untranslated\_block}（附录证明的数学密集句保留英文）、\texttt{font\_size\_drift}（字号缩放）、\texttt{display\_formula\_misaligned}（行间公式横向偏移 $\le$25pt）。
\end{itemize}
\end{frame}
\end{document}
